\subsection{5.1.4. RL 中的可控思考强度}\label{controllable-reasoning-effort-in-rl}

在模型架构与硬件取得进展的同时，输出 token 的数量同样是部署成本
的另一关键决定因素，也因此直接影响现实应用中的成本--质量权衡。
为此，我们在强化学习（RL）过程中引入一个标量强度档位 $b$ 作为显式的
条件信号。该机制同时应用于单轮推理与多轮智能体（agentic）任务。
具体而言，我们在系统提示词前附加如下指令：

思考强度（Reasoning Effort）: \{effort\} （范围 1--100；数值越大表示要求
更彻底的推理）

这里，\(b \in \{ 1 , \dots , 1 0 0 \}\) 表示请求的努力
层级。

对于每个训练提示 $x$，我们在每个强度档位 \(b \in { \mathcal { B } } :\) 上采样 \(M _ { b }\) 个响应：

\[
z _ { b , j } \sim \pi _ { \boldsymbol { \theta } } ( \cdot \mid x , b ) , \qquad b \in \mathcal { B } , \quad j = 1 , \ldots , M _ { b } .\tag{8}
\]

这里，$j$ 索引在强度档位 $b$ 上采样的响应。共享相同 \(( x , b )\) 的响应构成一个子组，组内奖励经过均值中心化以计算组间相对优势。因此，来自
不同强度档位的响应不会被直接比较。取而代之的是，通过让奖励的
长度分量依赖于 $b$，在每个子组内部诱导出随强度档位变化的行为。具体而言，我们将
长度惩罚项 \(r _ { b , j } ^ { \mathrm { l e n } }\) 加入
响应 \(z _ { b , j } \mathrm { : }\) 的奖励

\[
r _ { b , j } ^ { \mathrm { l e n } } = - \operatorname* { m i n } \left\{ C _ { \mathrm { m a x } } , \ k ( b ) \frac { \ell _ { b , j } } { L _ { \mathrm { n o r m } } } \right\} ,\tag{9}
\]

其中 \(\ell _ { b , j }\) 表示推理 token 的数量，
\(L _ { \mathrm { n o r m } }\) 是参考长度，
\(C _ { \mathrm { m a x } }\) 限制单条轨迹可承受的最大扣除。token 惩罚系数随
请求的强度档位呈指数衰减：

\[
k ( b ) = k _ { 0 } \exp \left( - \frac { b - b _ { \mathrm { m i n } } } { \tau } \right) , \qquad \tau = \lambda \overline { { \Delta b } } ,\tag{10}
\]

其中 \(k _ { 0 }\) 是不同强度档位下的基础惩罚系数，
\(b _ { m i n }\) 是 $\mathcal{B}$ 的最小值，$\overline { { \Delta b } }$ 是训练强度档位之间的平均
间隔，而 $\lambda$ 控制惩罚衰减的速率。将 $b$ 增大 $\tau$ 会使惩罚系数乘以
\(e ^ { - 1 }\) 。参数 \(k _ { 0 }\) 控制整体偏向更短推理的力度，
而较小的 $\tau$ 会使惩罚衰减得更快，倾向于在
强度档位之间产生更大的行为差异。附录 C 为 $k(b)$ 的指数参数化提供了边际效用层面的
动机解释。



在部署阶段，标量 $b$ 提供了对模型思考强度的灵活控制接口。
通过改变 $b$，同一个模型 checkpoint 可以沿已学习的成本--质量前沿在不同的运行区间之间移动，以适应不同的延迟、
token 预算和求解质量要求。尽管训练时
只用到一个有限的强度档位集合，部署时仍可使用中间取值以引出
插值式的推理行为，从而为测试期的资源分配提供一种
细粒度且高效的机制。

在我们于 2026 年 9 月上线的生产部署中，公共 API
暴露了三个预设的思考强度档位---max、high、low（高、中、低）---它们
直接映射到这一标量接口上。如表 2 所述，这三个档位分别对应强度值 $b = 100$、$b = 75$、$b = 50$，因此 API 用户可以在已学习的成本--质量前沿上选择运行点，而无需改动模型权重或
解码配置。

表 2 \textbar{} 公共 API 思考强度档位与底层标量强度值 $b$ 之间的对应关系。

{\def\LTcaptype{none} % do not increment counter
\begin{longtable}[]{@{}|l|l|@{}}
\toprule\noalign{}
\endhead
\bottomrule\noalign{}
\endlastfoot
\hline
API 档位 & 强度值 $b$ \\
\hline
max & 100 \\
\hline
high & 75 \\
\hline
low & 50 \\
\hline
\end{longtable}
}

\subsection{5.2. 异步后训练基础设施}\label{asynchronous-post-training-infrastructure}

LLM 强化学习（RL）rollout 阶段的长尾问题一直是训练效率的
主要瓶颈。为解决这一问题，我们扩展后训练基础设施以支持异步
样本生成（Zeng et al., 2026; Team et al., 2026b），通过
维持足够高的并发度显著缓解 rollout 阶段的长尾问题。目前，异步
训练已在我们几乎所有的 RL 和 OPD 任务中启用，rollout
效率也获得了大幅提升。

\subsection{5.2.1. 总体工作流}\label{overall-workflow}

我们将 rollout 与训练共同部署在同一批物理设备上，并对二者的执行进行
时分复用，从而无需手动调整两个阶段之间的资源分配。每个任务都会指定在途样本数量的上界，
系统在整个 rollout 阶段维持该
上界。

我们评估了三种用于维持目标 rollout 并发度的调度粒度。
最终采用的是样本级调度：一旦
新完成的样本数达到分配给下一个提示的 GRPO 组大小，
无论这些完成样本来自哪些组，我们都立即调度该提示。这有助于在整个训练过程中维持稳定的 rollout
并发度。在确定这一调度
策略之前，我们试验了另外两种调度粒度。
第一种做法是在开始时额外调度几个批次，并在每个训练迭代后补充一整批；
然而这导致训练指标出现剧烈振荡，表明
批级粒度过于粗糙。第二种做法
改为提示级调度，即在一个 GRPO 组完成后调度新的提示，
但发现它会轻易卡在 GRPO 组内的长尾样本上，难以平稳地
维持目标 rollout 并发度。

一旦积累了足够的训练样本，训练便会抢占正在进行的
rollout。训练期间，我们使用拼接式路由回放（concatenated routing-replay）：对于
跨越多个 checkpoint 的样本，我们拼接每个 rollout 段产生的专家
路由，而不是丢弃这些信息并利用新 checkpoint
重新计算路由。



\subsection{5.2.2. 缓解长度偏差与离策略影响}\label{mitigating-length-bias-and-off-policy-effects}

异步生成在有效提升 rollout 效率的同时，
也引入了两个可能损害训练质量的副作用。其一，它
会造成长度分布偏差，尤其是在训练早期
阶段，因为较短的序列往往先完成，从而
主导最初的训练批次。其二，它不可避免地产生
离策略（off-policy）样本，即 token 部分或全部由较早的 checkpoint 生成的样本。这两类问题需要采取不同的
处理策略。

针对长度偏差，我们采用两种机制。第一，
调度器可以按数据集限制并发度，这有助于
调节每个数据集在稳态训练批次中的占比，通过控制新到样本的来源
间接缓解长度偏斜。第二，我们支持丢弃提前返回的短
样本，以平滑过渡到稳态长度
分布，并防止模型对过短的
序列过拟合。

针对离策略问题，我们实现了另外两个机制。第一，
通过调整控制样本调度与训练样本等待条件的逻辑，我们可将最大离策略
比例限制在某个上界内，确保训练数据不会过度偏离
当前模型。第二，在训练过程中，我们加入一种损失掩码方案，
剔除陈旧度过高的 token 的贡献，从而减轻陈旧样本对梯度
更新的不利影响。